\makeatletter
\def\input@path{{./}{../}{../../}{../../../}{../../../../}{preamble/}{../preamble/}{../../preamble/}{../../../preamble/}{../../../../book2026-2027/preamble/}}
\makeatother
\input{preamble_exams.tex}
\input{graphs/highlights.tex}
% \GTenHighlightsfalse % Uncomment to disable instructional highlights.
% Population/model, sample evidence, design threshold, and result roles.
\newcommand{\Pop}[1]{\HighlightStroke{GTenHLParameter}{#1}}
\newcommand{\Obs}[1]{\HighlightStroke{GTenHLValue}{#1}}
\newcommand{\Design}[1]{\HighlightStroke{GTenHLThird}{#1}}
\newcommand{\Result}[1]{\HighlightStroke{GTenHLOutcome}{#1}}

\title{Grade 11 Global Economics\\Term 3 --- Lesson 1\\Home Practice C8: Independent two-sample t-tests}
\author{}\date{}
\hypersetup{pdftitle={Grade 11 Global Economics - Term 3 - Lesson 1 - C8},pdfsubject={Independent two-sample t-tests}}
\begin{document}
\maketitle
\setcounter{tocdepth}{1}
\setcounter{secnumdepth}{2}
\phantomsection\label{toc}
\tableofcontents
\section*{Evaluated criteria}
\begin{description}
\item[\textbf{C8}] Conducts independent two-sample t-tests to compare the means of two populations, verifying the required conditions and interpreting the results in economic and financial contexts.
\end{description}
\textbf{How to practise.} Check whether the samples are independent or matched before computing. Use Welch's test and keep the subtraction order consistent. Attempt each task before opening its solution link. Use the return arrows to revisit the problem, group, or contents. All datasets and scenarios in this activity are illustrative.
\clearpage
\section{Introduction}\label{sec:introduction}
\subsection*{Highlight legend}
\phantomsection\label{tab:highlight-legend}
\begin{tabularx}{\linewidth}{@{}>{\raggedright\arraybackslash}p{.22\linewidth}>{\raggedright\arraybackslash}p{.25\linewidth}>{\raggedright\arraybackslash}X@{}}
\toprule
\textbf{Symbol or notation}&\textbf{Concept}&\textbf{Brief meaning}\\\midrule
$\Pop{\mu_1-\mu_2}$ & \Pop{Population contrast} & Difference of population means. \\
$\Obs{\bar x_1-\bar x_2}$ & \Obs{Observed contrast} & Estimated mean difference. \\
$\Obs{s_1},\Obs{s_2},\Obs{n_1},\Obs{n_2}$ & \Obs{Group summaries} & Separate spreads and sizes. \\
$\Obs{t},\Design{\nu}$ & \Design{Welch reference} & Statistic and degrees of freedom. \\
\bottomrule
\end{tabularx}
\subsection{Introductory example 1: Two independent cities}\label{ex:intro-one}
Random samples of different households in cities A and B give means 105 and 95 USD weekly spending. Define population contrast $\Pop{\mu_1-\mu_2}$ as A minus B, estimated by $\Obs{\bar x_1-\bar x_2}=\Obs{10}$ USD. No household appears in both samples.
\subsection{Introductory example 2: Unequal group precision}\label{ex:intro-two}
Independent samples of businesses have $n_1=9$, $\bar x_1=110$, $s_1=12$ and $n_2=16$, $\bar x_2=100$, $s_2=8$, all monetary summaries in USD. Both populations are approximately normal. Welch's model retains each group's variance contribution; it does not require them to be equal.
\section{Notation Reference and Practice Procedure}\label{sec:notation-and-procedure}
\subsection{Notation reference}\label{tab:notation-reference}
\begin{tabularx}{\linewidth}{@{}>{\raggedright\arraybackslash}p{.16\linewidth}>{\raggedright\arraybackslash}p{.18\linewidth}>{\raggedright\arraybackslash}p{.28\linewidth}>{\raggedright\arraybackslash}X@{}}
\toprule
\textbf{Symbol or notation}&\textbf{Concept}&\textbf{Definition}&\textbf{Meaning in the introductory examples}\\\midrule
$\Pop{\mu_1-\mu_2}$ & \Pop{Population contrast} & The subtraction order determines the sign and the alternative; units are those of the outcome. & \textbullet\ \hyperref[ex:intro-one]{Example 1}: city A minus city B mean spending. \\
$\Obs{\bar x_1-\bar x_2}$ & \Obs{Observed contrast} & Each sample mean summarizes a separate group; it is not a mean of matched differences. & \textbullet\ \hyperref[ex:intro-one]{Example 1}: $105-95=10$ USD. \\
$\Obs{s_1},\Obs{s_2},\Obs{n_1},\Obs{n_2}$ & \Obs{Group summaries} & Welch uses $SE=\sqrt{s_1^2/n_1+s_2^2/n_2}$ without assuming equal population variances. & \textbullet\ \hyperref[ex:intro-two]{Example 2}: different sample spreads and sample sizes. \\
$\Obs{t},\Design{\nu}$ & \Design{Welch reference} & $t=[(\bar x_1-\bar x_2)-\Delta_0]/SE$; $\nu$ is the Welch--Satterthwaite degrees of freedom, often non-integer. & \textbullet\ \hyperref[ex:intro-two]{Example 2}: compare a 10 USD estimated gap with zero. \\
\bottomrule
\end{tabularx}
\subsection{Practice procedure}
\begin{enumerate}
\item Define both populations and contrast order, benchmark $\Delta_0$, alternative direction, and planned level.
\item Check random sampling, independence within and between groups, and that each person or firm occurs only in the intended group. Matched observations require a paired analysis.
\item For small groups require approximately normal populations without serious outliers; larger samples do not excuse severe dependence or heavy tails.
\item Compute $SE$, Welch $t$, and $\nu$ using separate sample variances. Use the t reference with that degrees of freedom.
\item Compare statistic to the planned critical region or p-value to $\alpha$. Interpret the signed effect, sampling scope, and economic importance.
\end{enumerate}
\section{General Worked Problems}\label{sec:general-worked}
\SectionProblemsTableOfContents{\SectionProblemLink{prob:equal}{City spending comparison} \\
\SectionProblemLink{prob:unequal}{Unequal variances and sample sizes}}
\subsection{City spending comparison}\label{prob:equal}
Independent simple random samples from large city populations give $n_1=n_2=10$, $\bar x_1=105$ USD, $\bar x_2=95$ USD, and $s_1=s_2=10$ USD. Both populations are approximately normal without serious outliers. Test any mean difference at $\alpha=0.05$. Supplied for $df=18$: two-sided critical values $\pm2.101$; $P(|T_{18}|\ge2.236)\approx0.0382$.
\QuestionsTasks[sec:general-worked]{{8}/{sol:equal}/{Conduct Welch's test and explain whether equal sample SDs require pooling.}}
\subsubsection*{Worked solution}\label{sol:equal}
Let $\mu_1$ and $\mu_2$ be the two population mean weekly expenditures. $H_0:\mu_1-\mu_2=0$, $H_a:\mu_1-\mu_2\ne0$. Independent random samples and normal populations justify the small-sample model. Welch gives
\[
SE=\sqrt{10^2/10+10^2/10}=\sqrt{20},\qquad
t=\frac{\Obs{105}-\Obs{95}-\Pop{0}}{\sqrt{20}}\approx2.236.
\]
For $a=s_1^2/n_1$ and $b=s_2^2/n_2$, the degrees of freedom are
\[
\nu=\frac{(a+b)^2}{a^2/(n_1-1)+b^2/(n_2-1)}
=\frac{20^2}{10^2/9+10^2/9}=18.
\]
Since $2.236>2.101$ and p-value $0.0382<0.05$, \Result{reject $H_0$}. Evidence supports different population means, with A estimated 10 USD higher. Equal observed SDs do not establish equal population variances; Welch remains appropriate without a pooled assumption. Random sampling supports population comparison, not a causal effect of city residence.
\secInicio[prob:equal][sec:general-worked]
\subsection{Unequal variances and sample sizes}\label{prob:unequal}
Two independent random samples of different businesses have $n_1=9$, $\bar x_1=110$, $s_1=12$ and $n_2=16$, $\bar x_2=100$, $s_2=8$, in USD. Both populations are approximately normal without serious outliers. The preplanned question is whether population 1 has a higher mean at $\alpha=0.05$. Supplied: Welch $\nu\approx12.10$, right-tail critical approximately 1.781; at $t=2.236$, the right-tail p-value lies between 0.01 and 0.025.
\QuestionsTasks[sec:general-worked]{{8}/{sol:unequal}/{Conduct the test and explain the contributions to uncertainty.}}
\subsubsection*{Worked solution}\label{sol:unequal}
Use $H_0:\mu_1-\mu_2=0$, $H_a:\mu_1-\mu_2>0$. Independent sampling and normal populations are given; unequal population variances need not be ruled out. Here
\[
a=12^2/9=16,\quad b=8^2/16=4,\quad SE=\sqrt{20},
\quad t=(110-100)/\sqrt{20}\approx2.236,
\]
\[
\nu=\frac{(16+4)^2}{16^2/8+4^2/15}\approx12.10.
\]
The first group contributes more to the squared standard error because its spread is larger and its sample smaller. Since $2.236>1.781$ and the supplied p-value bound is entirely below 0.05, \Result{reject $H_0$}. Evidence supports a higher population mean for group 1; the estimated gap is 10 USD. Fractional degrees of freedom are valid. They should not be replaced by $n_1+n_2-2$ unless an equal-variance pooled model is justified.
\secInicio[prob:unequal][sec:general-worked]
\section{Advanced Problems}\label{sec:advanced}
\SectionProblemsTableOfContents{\SectionProblemLink{prob:design}{Similar-looking tables, different designs}}
\subsection{Similar-looking tables, different designs}\label{prob:design}
Study A measures loan-processing cost before and after software installation at the same 15 branches. Study B samples 15 different branches with the software and 18 other branches without it. In B, all branches belong to one of three banking chains and branches in a chain share management shocks. The manager proposes Welch's test for both studies and says a small p-value would prove software causes savings.
\QuestionsTasks[sec:advanced]{{8}/{sol:design}/{Select the appropriate structure and identify remaining conditions and interpretation limits.}}
\subsubsection*{Worked solution}\label{sol:design}
A is paired: define $d=\text{before}-\text{after}$ and study population $\mu_d$ using 15 differences, provided pairs are independent and differences satisfy shape conditions. Welch treats groups as independent and would discard the match. B has separate groups, so the intended contrast is mean cost with software minus mean cost without software; Welch is the default only if within-group and between-group independence are justified. Shared chain shocks create clustering, which ordinary Welch standard errors do not handle. Obtain a suitable cluster-based design or dependence-aware method; do not report a justified Welch decision from the description. Neither observational design alone proves causation. Compare estimated savings with installation and maintenance costs even if a valid future analysis finds statistical significance.
\secInicio[prob:design][sec:advanced]
\end{document}

